% !TEX program = xelatex
% 공통수학2 · Ⅱ. 집합과 명제 · 1. 집합 · 04 여집합과 차집합 · 2/2차시
% 학생용 활동지 (답안 없음)
\documentclass[11pt,a4paper]{article}
\usepackage{kotex}
\usepackage{iftex}
\ifPDFTeX\else
  \setmainhangulfont{Noto Serif CJK KR}
  \setsanshangulfont{Noto Sans CJK KR}
\fi
\usepackage[margin=2.1cm]{geometry}
\usepackage{parskip}
\usepackage{enumitem}
\usepackage{array}
\usepackage{tabularx}
\usepackage{xcolor}
\usepackage{amssymb}
\usepackage{tikz}
\usepackage[most]{tcolorbox}
\usepackage{fancyhdr}
\usepackage{titlesec}

\definecolor{goldc}{HTML}{B07C22}
\definecolor{goldstrong}{HTML}{8A5F16}
\definecolor{indigoc}{HTML}{2B3B57}
\definecolor{indigosoft}{HTML}{6E82A6}
\definecolor{linec}{HTML}{D6DACB}
\definecolor{surface2}{HTML}{E4E8DC}
\definecolor{writeline}{HTML}{C7CCBA}
\definecolor{inksoft}{HTML}{52604F}

\titleformat{\section}{\normalfont\sffamily\bfseries\large\color{indigoc}}{\thesection}{0.6em}{}
\titlespacing{\section}{0pt}{14pt}{6pt}
\newtcolorbox{goalbox}{enhanced, breakable, colback=white, colframe=goldc, boxrule=0.5pt, leftrule=3pt, arc=2pt, left=10pt, right=10pt, top=8pt, bottom=8pt}
\newtcolorbox{sheetbox}{enhanced, breakable, colback=white, colframe=linec, boxrule=0.6pt, arc=3pt, left=11pt, right=11pt, top=9pt, bottom=9pt}
\newcommand{\writespace}[1]{%
  \par\nobreak\vspace{3pt}
  \foreach \n in {1,...,#1} {%
    \noindent\textcolor{writeline}{\rule{\linewidth}{0.4pt}}\par\vspace{0.62cm}%
  }
}
\newcommand{\fillblank}[1]{\makebox[#1]{\hrulefill}}
\pagestyle{fancy}\fancyhf{}\renewcommand{\headrulewidth}{0pt}
\fancyfoot[L]{\footnotesize\color{inksoft} 공통수학2 · 2026학년도 2학기 · 지평선고등학교}
\fancyfoot[R]{\footnotesize\color{inksoft} 다음 시간 --- 중단원 마무리(집합)}

\begin{document}

\noindent
\begin{minipage}[b]{0.62\linewidth}
  {\sffamily\footnotesize\color{indigosoft} 공통수학2 · Ⅱ. 집합과 명제 · 1. 집합}\\[2pt]
  {\sffamily\bfseries\Large 04 여집합과 차집합 \textperiodcentered\ 24차시}
\end{minipage}\hfill
\begin{minipage}[b]{0.36\linewidth}
  \raggedleft\small\color{inksoft}
  반\ \fillblank{1.6cm} \quad 번호\ \fillblank{1.6cm}\\[4pt]
  이름\ \fillblank{3.6cm}
\end{minipage}

\vspace{4pt}
\noindent{\color{goldc}\rule{\linewidth}{2.2pt}}
\vspace{10pt}

\begin{goalbox}
\textbf{오늘의 목표}\\
여집합·차집합의 성질과 드모르간의 법칙을 이해하고, 이를 이용해 원소의 개수를 구할 수 있다.
\vspace{4pt}
\small\color{inksoft}
$\square$ 자신 있음 \quad $\square$ 조금 헷갈림 \quad $\square$ 처음 봄 \hfill (수업 전 나의 상태에 표시)
\end{goalbox}

\section{개념 정리 --- 여집합·차집합의 성질과 드모르간의 법칙}

\begin{sheetbox}
전체집합 $U$의 두 부분집합 $A$, $B$에 대하여

\begin{itemize}[leftmargin=1.6em, itemsep=3pt]
  \item $U^C=\ \fillblank{2cm}$, \quad $\varnothing^C=\ \fillblank{2cm}$
  \item $(A^C)^C=\ \fillblank{2cm}$
  \item $A\cup A^C=\ \fillblank{2cm}$, \quad $A\cap A^C=\ \fillblank{2cm}$
  \item $A-B=A\cap\ \fillblank{2cm}$
\end{itemize}

\vspace{10pt}
\textbf{드모르간의 법칙}
\[
(A\cup B)^C = \fillblank{3.2cm} \qquad\qquad (A\cap B)^C = \fillblank{3.2cm}
\]

{\footnotesize\color{inksoft} 두 식에서 $\cup$과 $\cap$이 어떻게 바뀌는지 한 문장으로 적어 보자.}
\writespace{1}
\end{sheetbox}

\section{연습 문제}

\begin{sheetbox}
\textbf{1.} 전체집합 $U$의 두 부분집합 $A$, $B$에 대하여 $n(U)=50$, $n(A)=30$, $n(B)=21$, $n(A\cap B)=10$일 때, $n(A^C\cap B^C)$를 구하시오.
\writespace{4}

\vspace{6pt}
\textbf{2.} 전체집합 $U$의 두 부분집합 $A$, $B$에 대하여 $n(U)=40$, $n(A)=21$, $n(B)=17$, $n(A\cup B)=34$일 때, $n(A^C\cup B^C)$를 구하시오.
\writespace{4}
\end{sheetbox}

\section{실생활 탐구}

\begin{sheetbox}
어느 학급의 학생 25명 중 직업 체험과 대학 탐방을 모두 신청한 학생은 6명, 둘 다 신청하지 않은 학생은 2명이다. 직업 체험을 신청한 학생이 대학 탐방을 신청한 학생보다 5명 더 많을 때, 다음을 구하시오.

\begin{enumerate}[leftmargin=1.4em, itemsep=2pt]
  \item 직업 체험 또는 대학 탐방을 신청한 학생 수
  \item 직업 체험을 신청한 학생 수
\end{enumerate}
\writespace{6}
\end{sheetbox}

\section{사고의 가시화 --- Connect · Extend · Challenge}

\noindent{\footnotesize\color{inksoft} `여집합과 차집합' 전체(23$\sim$24차시)를 떠올리며 아래 세 칸을 채워 보자.}\\[6pt]
\begin{tabularx}{\linewidth}{@{}XXX@{}}
\textbf{\footnotesize\color{indigoc} Connect --- 무엇과 연결되나?} &
\textbf{\footnotesize\color{indigoc} Extend --- 어디까지 확장됐나?} &
\textbf{\footnotesize\color{indigoc} Challenge --- 아직 궁금한 것은?} \\[4pt]
\end{tabularx}
\noindent
\begin{minipage}[t]{0.32\linewidth}\writespace{3}\end{minipage}\hfill
\begin{minipage}[t]{0.32\linewidth}\writespace{3}\end{minipage}\hfill
\begin{minipage}[t]{0.32\linewidth}\writespace{3}\end{minipage}

\section{마무리 성찰}

\begin{sheetbox}
\begin{minipage}[t]{0.48\linewidth}
{\footnotesize\color{inksoft} `집합' 대단원을 마치며 한 문장으로}
\writespace{2}
\end{minipage}\hfill
\begin{minipage}[t]{0.48\linewidth}
{\footnotesize\color{inksoft} 감사 일기 / 유념 한 줄}
\writespace{2}
\end{minipage}
\end{sheetbox}

\end{document}
